\chapter{Configuration reference}
\label{app:configref}

Every environment variable and TOML key: name, whether required, format, example,
default, and effect if missing.

\begin{center}
\begin{tabular}{@{}>{\ttfamily}l l p{6.5cm}@{}}
\toprule
\normalfont\textbf{Variable} & \textbf{Req.} & \textbf{Effect} \\
\midrule
SMARTME\_CLIENT\_ID     & yes & smart-me OAuth client-credentials ID. Missing: \prog refuses to start. \\
SMARTME\_CLIENT\_SECRET & yes & smart-me OAuth client-credentials secret. Missing: \prog refuses to start. \\
SMARTME\_BASIC\_USER    & --- & HTTP Basic fallback user. \textbf{Not read by \prog}\textsuperscript{*}. \\
SMARTME\_BASIC\_PASSWORD& --- & HTTP Basic fallback password. \textbf{Not read by \prog}\textsuperscript{*}. \\
\addlinespace
SMARTME\_LOG\_DIR       & no  & Directory for the log file. Unset: no log file, console only. Set: \code{smartme\_mqtt.<date>.log}, rotated daily. \\
SMARTME\_LOG\_KEEP      & no  & Rotated files kept. Default 7. Not a positive number: refused on stderr, keeps 7. \\
RUST\_LOG               & no  & \code{tracing} filter, console and file alike. Default \code{info}. \\
\bottomrule
\end{tabular}
\end{center}

\footnotesize\textsuperscript{*}\textbf{The Basic fallback is not reachable from the
environment.} \code{smart-me-client} implements it (ADR~0009), but the binary constructs
OAuth client-credentials unconditionally and never builds Basic credentials, so setting
these two variables has no effect --- and because the OAuth pair is required, omitting it in
their favour makes \prog refuse to start. Wiring it up, or dropping it from the design, is an
open decision.\normalsize

\stub{the full table (broker, meter identity, Sparkplug identity, state directory, poll
interval, bind, mapping) as those settings are documented here.}
